%----------------------------------------------------------------------------------------
% RAPPORTEREN. Alle opgaven zijn eigen werk.
%
% Opbouw: Start, leerstof (theorie, werkwijze of uitleg, voorbeeld, korte oefening), en op een
% nieuwe pagina de Opdrachten (door elkaar, stijgende moeilijkheid), Koppel en Test jezelf.
% Labels van dit hoofdstuk beginnen met h3-.
%----------------------------------------------------------------------------------------

\mychapter{3}{Rapporteren}

\lkr{Leerplan B+S'': LPD 55 (onderzoekscompetentie), het deel over rapporteren en reflecteren. Wat hier staat over de vorm van een verslag is een voorstel; volg de afspraken van de school als die anders zijn (bijvoorbeeld voor bronvermelding of het gebruik van AI-hulpmiddelen). Laat leerlingen elkaars tekst nalezen met de checklist: dat leert meer dan de leerkracht alleen.}

Een onderzoek dat niemand begrijpt, heeft weinig zin. Een goed verslag laat de lezer je redenering volgen, en overtuigt met argumenten, niet met veel tekst. Wiskundig schrijven is een vaardigheid apart: precies, kort, en met formules die in zinnen staan. Dit hoofdstuk gaat over de opbouw van een verslag, over wiskunde en getallen opschrijven, en over figuren, tabellen en presentaties.

\begin{oefeningen}[Start]
\oef{h3-start-1}

\oef{h3-verken}
\end{oefeningen}

%========================================================================================
\section{Een verslag opbouwen}

\subsection{De onderdelen van een verslag}

Een verslag is geen dagboek van alles wat je deed, in de volgorde waarin het gebeurde. Je ordent het zo dat een lezer snel vindt wat hij zoekt: wat was de vraag, hoe heb je gewerkt, wat vond je, en wat is het antwoord? Daarvoor bestaat een vaste opbouw.

\begin{theorie}
Een onderzoeksverslag bestaat uit:
\begin{itemize}
\item \textbf{\textcolor{accent}{inleiding}}: de onderzoeksvraag, waarom ze interessant is, en de deelvragen;
\item \textbf{\textcolor{accent}{methode}}: hoe je gewerkt hebt (metingen, modellen, redeneringen), zo dat iemand anders het kan herhalen;
\item \textbf{\textcolor{accent}{resultaten}}: tabellen, figuren en berekeningen, met uitleg;
\item \textbf{\textcolor{accent}{besluit}}: het antwoord op de onderzoeksvraag, en de grenzen ervan;
\item \textbf{\textcolor{accent}{reflectie}}: wat ging goed, wat zou je anders doen, welke nieuwe vragen;
\item \textbf{\textcolor{accent}{bronnen}}: alles wat je gebruikt hebt.
\end{itemize}
\end{theorie}

\voorbeeld{}{Schrijf voor het onderzoek naar de slinger uit hoofdstuk 2 bij elk onderdeel van het verslag één of twee zinnen.\label{vb:h3-opbouw}}{%
\uitleg{Per onderdeel vragen we ons af wat de lezer daar wil weten. De zinnen zelf staan cursief.}
{\renewcommand{\arraystretch}{1.2}%
\begin{tabular}{@{}>{\RaggedRight\arraybackslash}p{2.2cm}@{\hspace{4mm}}>{\RaggedRight\itshape\arraybackslash}p{12.3cm}@{}}
inleiding & Hoe hangt de periode $T$ van een slinger af van zijn lengte $L$? Slingers regelden eeuwenlang klokken, dus dat verband bepaalt hoe je een klok bijstelt. \\
methode & Voor vijf lengtes van $0{,}2$ m tot $1{,}0$ m maten we drie keer de tijd voor tien periodes, telkens met hetzelfde gewichtje en een kleine beginuitwijking. \\
resultaten & Tabel 1 en figuur 1 tonen $T^2$ tegen $L$. De kleinste-kwadratenrechte is $T^2 = 4{,}025\,L - 0{,}003$, met $r \approx 1{,}000$. \\
besluit & Voor lengtes tussen $0{,}2$ m en $1{,}0$ m is $T^2$ evenredig met $L$: $T \approx 2{,}01\sqrt{L}$, met $T$ in s en $L$ in m. \\
reflectie & Door tien periodes te meten, speelde onze reactietijd nauwelijks mee. Een nieuwe vraag: blijft het verband gelden bij een grote uitwijking? \\
bronnen & Het natuurkundeboek waaruit we de formule $T = 2\pi\sqrt{L/g}$ haalden, met auteur, titel, uitgever en jaar.
\end{tabular}}}

\begin{oefeningen}
\oef{h3-onderdelen}
\end{oefeningen}

\subsection{Besluit en reflectie}

Het besluit moet passen bij de vraag uit de inleiding: lees ze op het einde nog eens na elkaar. Vraagt de inleiding \textit{hoeveel}, dan staat in het besluit een getal. Vraagt ze \textit{hoe hangt \ldots\ af van}, dan staat er een verband, met de grenzen waarbinnen het getest is. Een besluit als \textit{we hebben veel geleerd} is geen antwoord.

De reflectie gaat niet over het onderwerp, maar over het onderzoek zelf: hoe verliep het, wat zou je anders doen, en welke nieuwe vragen kwamen er op? Zo hoort \textit{onze metingen bij korte slingers waren minder nauwkeurig} in de reflectie, en niet in het besluit.

\voorbeeld{}{Onderzoeksvraag: \textit{Hoeveel diagonalen heeft een veelhoek met $n$ hoekpunten?} Beoordeel twee besluiten, en schrijf een reflectie.\label{vb:h3-besluit}}{%
\uitleg{We lezen elk besluit samen met de vraag: geeft het een antwoord op \textit{hoeveel}, voor elke $n$?}
\begin{tabular}{@{}>{\RaggedRight\arraybackslash}p{6.4cm}@{\hspace{4mm}}>{\RaggedRight\small\itshape\arraybackslash}p{8.3cm}@{}}
\textit{We vonden een mooie formule die altijd klopt.} & Geen antwoord: welke formule? En waarom klopt ze altijd? \\[.6ex]
\textit{Een veelhoek met $n$ hoekpunten heeft $\frac{n(n - 3)}{2}$ diagonalen. We bewezen dat door per hoekpunt te tellen.} & Goed: een antwoord voor elke $n$, met de reden waarom het zeker is. \\[.6ex]
\textit{We vonden de formule pas toen we de tabel vergeleken met de handdrukken. Nieuwe vraag: hoeveel snijpunten hebben de diagonalen binnen de veelhoek?} & Een reflectie: hoe het onderzoek verliep, en een nieuwe vraag.
\end{tabular}}

\begin{oefeningen}
\oef{h3-besluit}
\end{oefeningen}

\subsection{Bronnen en eerlijkheid}

Alles wat niet van jezelf komt, vermeld je: een formule uit een boek, gegevens van een website, een idee van een klasgenoot, hulp van een AI-hulpmiddel. Dat is eerlijk tegenover anderen, en het laat de lezer toe je bronnen na te kijken. Algemeen bekende leerstof, zoals de stelling van Pythagoras, hoef je niet te vermelden.

\begin{theorie}
\textbf{\textcolor{accent}{Plagiaat}} is werk van iemand anders gebruiken zonder de bron te vermelden, ook als je de tekst een beetje herschrijft.

\vspace{1mm}
Bij een \textbf{\textcolor{accent}{bron}} vermeld je minstens: \textbf{wie} (auteur of organisatie), \textbf{wat} (titel), \textbf{waar} (uitgever of webadres) en \textbf{wanneer} (jaar, of bij een website de datum waarop je ze raadpleegde).
\end{theorie}

\voorbeeld{}{Je verslag gebruikt een tabel met bevolkingscijfers van de website van een statistiekbureau. Een AI-hulpmiddel las je tekst na op spelling. Hoe vermeld je dat?\label{vb:h3-bron}}{%
\uitleg{Voor de website zoeken we de vier gegevens: wie, wat, waar en wanneer. Bij het AI-hulpmiddel zeggen we welk het was, en waarvoor we het gebruikten.}
\begin{stappenvrij}{5.4cm}{9.3cm}
\text{wie: het statistiekbureau} & De organisatie, want een auteur staat er niet bij. \\
\text{wat: de titel van de tabel} & Zo vindt de lezer precies dezelfde cijfers. \\
\text{waar: het webadres} & Het volledige adres van de pagina, niet alleen de startpagina. \\
\text{wanneer: de raadpleegdatum} & Een website verandert: de cijfers kunnen later anders zijn. \\
\text{AI-hulpmiddel: in de methode} & \textit{De tekst werd op spelling nagelezen met [naam van het hulpmiddel]. De berekeningen en besluiten zijn van ons.}
\end{stappenvrij}
Volg voor de precieze vorm de afspraken van je school.}

\begin{oefeningen}
\oef{h3-bronnen}
\end{oefeningen}

%========================================================================================
\section{Wiskunde opschrijven}

\subsection{Symbolen, formules en eenheden}

In een verslag ben je er niet bij om uit te leggen wat je bedoelde. Elke letter, elke formule en elk getal moet dus op zichzelf te begrijpen zijn. Een formule lees je trouwens als een zin: \textit{$g$ is gelijk aan vier pi kwadraat $L$ gedeeld door $T$ kwadraat.} Daarom staat ze ook in een zin, met een punt of een komma erachter.

\voorbeeld{}{Verbeter deze tekst uit een verslag: \textit{T = 2,006, dus g = 4pi² L/T² = 9,8083 want L = 1.}\label{vb:h3-tekst}}{%
\uitleg{Wat is $T$? Waar komt de formule vandaan? Hoe nauwkeurig is de meting? En wat is de eenheid?}
\begin{itemize}
\item \textit{Voor een slinger met lengte $L = 1{,}00$ m maten we een periode $T = 2{,}01$ s (gemiddelde van tien metingen).}
\item \textit{Uit $T = 2\pi\sqrt{L/g}$ volgt $g = \frac{4\pi^2 L}{T^2}$, en dus $g \approx 9{,}8$ m/s$^2$.}
\end{itemize}
De symbolen zijn uitgelegd voor ze gebruikt worden, de formule heeft een herkomst, en elk getal heeft een eenheid. Met een periode op twee decimalen is $g$ hooguit op één decimaal betrouwbaar: $9{,}8083$ suggereert een nauwkeurigheid die er niet is (zie \ref{sec:h3-rekenen}).}

\begin{werkwijze}
Definieer elk symbool voor je het gebruikt. Zet formules in zinnen: ook een formule eindigt met een punt of een komma. Schrijf bij elke stap waarom ze klopt. Geef getallen met een eenheid en met een zinvolle nauwkeurigheid: niet meer cijfers dan je metingen toelaten.
\end{werkwijze}

Ook de getallen in een model hebben een eenheid. In $T^2 = 4{,}03\,L$, met $T$ in s en $L$ in m, staat links s$^2$ en rechts $4{,}03$ maal m. Dus heeft $4{,}03$ de eenheid s$^2$/m. Zo heeft een richtingscoëfficiënt altijd de eenheid van $y$ gedeeld door die van $x$: bij een fiets die $45$ euro per jaar minder waard wordt, is dat euro per jaar.

\begin{oefeningen}
\oef{h3-schrijven}
\end{oefeningen}

\subsection{Nauwkeurigheid}

Een rekentoestel geeft tien cijfers, maar je metingen zijn niet zo nauwkeurig. Hoeveel cijfers mag je dan opschrijven? Om dat te beslissen, moet je eerst weten hoe nauwkeurig een gemeten getal is.

\begin{theorie}
Wijkt een gemeten of afgeronde waarde hoogstens $\Delta x$ af van de werkelijke waarde $x$, dan is $\Delta x$ de \textbf{\textcolor{accent}{absolute fout}} en $\formule{\frac{\Delta x}{x}}$ de \textbf{\textcolor{accent}{relatieve fout}}.

\vspace{1mm}
De \textbf{\textcolor{accent}{beduidende cijfers}} van een getal zijn alle cijfers vanaf het eerste cijfer dat geen $0$ is. Zo heeft $0{,}0450$ drie beduidende cijfers: $4$, $5$ en de laatste $0$.
\end{theorie}

De laatste $0$ in $0{,}0450$ telt mee: ze zegt dat je op een tienduizendste nauwkeurig gemeten hebt. Schrijf je $0{,}045$, dan beweer je minder. De absolute fout zegt hoeveel je ernaast kunt zitten, de relatieve fout hoe erg dat is: $1$ cm fout is veel op een potlood, en niets op een voetbalveld. Een relatieve fout geef je meestal in procent.

\voorbeeld{}{Een lengte is opgegeven als $12{,}40$ cm, gemeten op $0{,}01$ cm. Hoeveel beduidende cijfers heeft ze? Wat zijn de absolute en de relatieve fout?\label{vb:h3-fout}}{%
\uitleg{Gemeten op $0{,}01$ cm betekent dat de waarde is afgerond op honderdsten: de echte waarde ligt hoogstens een halve honderdste ernaast.}
\begin{stappenbreed}
\text{beduidende cijfers: } 1,\ 2,\ 4,\ 0 & Vier: de laatste $0$ zegt dat er op honderdsten gemeten is. \\
\Delta x = 0{,}005 \text{ cm} & De helft van de kleinste eenheid. \\
\frac{\Delta x}{x} = \frac{0{,}005}{12{,}40} \approx 0{,}0004 = 0{,}04\,\% & Een zeer nauwkeurige meting.
\end{stappenbreed}}

\begin{oefeningen}
\oef{h3-beduidend}
\end{oefeningen}

\subsection{Rekenen met afgeronde getallen}\label{sec:h3-rekenen}

Wat gebeurt er met de fouten als je met gemeten getallen rekent? Bij optellen en bij vermenigvuldigen gaat het anders, en dat verklaart de twee vuistregels voor afronden die je misschien al kent. We beginnen met een som.

\begin{theorie}
\begin{stelling}{Fout van een som}
Hebben $x$ en $y$ absolute fouten $\Delta x$ en $\Delta y$, dan hebben $x + y$ en $x - y$ een absolute fout van hoogstens \qquad $\formule{\Delta x + \Delta y}$
\end{stelling}

\begin{bewijs}
Stel dat je $x'$ en $y'$ gemeten hebt, met $|x' - x| \leqslant \Delta x$ en $|y' - y| \leqslant \Delta y$. Met de driehoeksongelijkheid $|u + v| \leqslant |u| + |v|$ is
\bewijsregel{\bigl|(x' + y') - (x + y)\bigr| = \bigl|(x' - x) + (y' - y)\bigr| \leqslant |x' - x| + |y' - y| \leqslant \Delta x + \Delta y .}
Voor het verschil schrijf je $(x' - y') - (x - y) = (x' - x) - (y' - y)$, en vind je dezelfde grens.
\end{bewijs}
\end{theorie}

Bij optellen en aftrekken tellen dus de absolute fouten op. De minst nauwkeurige term bepaalt dan bijna alles.

\voorbeeld{}{Je legt twee stukken touw achter elkaar: $2{,}35$ m, gemeten op $1$ cm, en $0{,}8$ m, geschat op $10$ cm. Hoe lang zijn ze samen?\label{vb:h3-som}}{%
\uitleg{We tellen de lengtes op, en ook hun absolute fouten.}
\begin{stappenbreed}
2{,}35 + 0{,}8 = 3{,}15 \text{ m} & De som van de gemeten waarden. \\
\Delta = 0{,}005 + 0{,}05 = 0{,}055 \text{ m} & De helft van $1$ cm plus de helft van $10$ cm. \\
\text{samen: } 3{,}2 \text{ m} & Op een fout van $5{,}5$ cm heeft een tweede decimaal geen zin.
\end{stappenbreed}}

Bij een product, zoals een oppervlakte, werkt het anders. Daar zijn de absolute fouten niet zo nuttig: een fout van $1$ cm op een zijde weegt zwaarder naarmate de andere zijde langer is. Wat wel eenvoudig optelt, zijn de relatieve fouten.

\begin{theorie}
\begin{stelling}{Fout van een product}
Hebben $x$ en $y$ relatieve fouten $\varepsilon$ en $\varphi$, dan heeft $xy$ een relatieve fout van hoogstens \qquad $\formule{\varepsilon + \varphi + \varepsilon\varphi \approx \varepsilon + \varphi}$
\end{stelling}

\begin{bewijs}
Schrijf de gemeten waarden als $x' = x(1 + e)$ en $y' = y(1 + f)$, met $|e| \leqslant \varepsilon$ en $|f| \leqslant \varphi$. Dan is
\bewijsregel{x'y' = xy(1 + e)(1 + f) = xy(1 + e + f + ef) ,}
zodat de relatieve fout van het product gelijk is aan
\bewijsregel{\frac{|x'y' - xy|}{|xy|} = |e + f + ef| \leqslant \varepsilon + \varphi + \varepsilon\varphi .}
\end{bewijs}
\end{theorie}

De term $\varepsilon\varphi$ is verwaarloosbaar: bij twee fouten van $1$\,\% is hij $0{,}01$\,\%. Voor een quotiënt vind je op dezelfde manier ongeveer $\varepsilon + \varphi$, en een kwadraat $T^2 = T \cdot T$ heeft de dubbele relatieve fout van $T$.

\voorbeeld{}{Een tafelblad is $1{,}20$ m lang en $0{,}75$ m breed, beide gemeten op $1$ cm. Hoe nauwkeurig ken je de oppervlakte?\label{vb:h3-nauwkeurig}}{%
\uitleg{Bij een product tellen de relatieve fouten op. We berekenen ze, en zetten de som om naar een absolute fout.}
\begin{stappenbreed}
\frac{0{,}005}{1{,}20} \approx 0{,}42\,\%,\quad \frac{0{,}005}{0{,}75} \approx 0{,}67\,\% & De relatieve fouten van lengte en breedte. \\
A = 1{,}20 \cdot 0{,}75 = 0{,}90 \text{ m}^2,\quad \text{fout} \approx 1{,}1\,\% & De relatieve fouten opgeteld. \\
\Delta A \approx 0{,}011 \cdot 0{,}90 \approx 0{,}01 \text{ m}^2 & De oppervlakte ligt tussen $0{,}89$ en $0{,}91$ m$^2$.
\end{stappenbreed}
Je schrijft dus $A \approx 0{,}90$ m$^2$: twee beduidende cijfers, net als de breedte $0{,}75$ m.}

\begin{werkwijze}
Bij een som of een verschil rond je af op de plaats van de minst nauwkeurige term: evenveel decimalen. Bij een product of een quotiënt geef je het resultaat met evenveel beduidende cijfers als de minst nauwkeurige factor. Rond pas af op het einde, niet in tussenstappen.
\end{werkwijze}

\begin{oefeningen}
\oef{h3-fouten}

\oef{h3-afronden}
\end{oefeningen}

%========================================================================================
\section{Figuren, tabellen en presentaties}

\subsection{Een goede figuur}

Een goede figuur kan zonder de tekst gelezen worden: de assen hebben een naam en een eenheid, de schaal is duidelijk, en er is een bijschrift dat zegt wat je moet zien. Metingen teken je als losse punten, eventueel met een model erbij. Verbind ze niet met lijnstukken: die suggereren waarden die je niet gemeten hebt.

\figuren{%
\begin{tikzpicture}[x=.7cm,y=.45cm,font=\scriptsize]
  \draw[->] (0,0) -- (5.6,0);
  \draw[->] (0,0) -- (0,6);
  \draw (1,2) -- (2,4) -- (3,5) -- (4,4) -- (5,5);
  \foreach \x/\y in {1/2,2/4,3/5,4/4,5/5}{\fill (\x,\y) circle (2pt);}
  \node[below] at (2.8,-.6) {\textit{niet goed}};
\end{tikzpicture}\hspace{1.2cm}
\begin{tikzpicture}[x=.7cm,y=.45cm,font=\scriptsize]
  \draw[black!14,very thin] (0,0) grid (5,6);
  \draw[->,grijs] (0,0) -- (5.6,0) node[right,text=black] {uren};
  \draw[->,grijs] (0,0) -- (0,6.6) node[above,text=black] {score (op 10)};
  \foreach \t in {1,...,5}{\node[below,inner sep=1.5pt] at (\t,0) {$\t$};}
  \foreach \t in {2,4,6}{\node[left,inner sep=1.5pt] at (0,\t) {$\t$};}
  \draw[accent,thick] (0,2.2) -- (5.5,5.5);
  \foreach \x/\y in {1/2,2/4,3/5,4/4,5/5}{\fill[accent] (\x,\y) circle (2pt);}
  \node[below] at (2.8,-.7) {\textit{beter}: assen met naam en schaal, met model};
\end{tikzpicture}}

\voorbeeld{}{Schrijf een bijschrift bij de rechtse figuur hierboven. Ze toont de score van vijf leerlingen op een test, tegen het aantal uren dat ze studeerden.\label{vb:h3-figuur}}{%
\uitleg{Een bijschrift zegt wat er staat, en wat de lezer moet opvallen.}
\textit{Figuur 2: testscore (op 10) tegen het aantal uren studie, voor vijf leerlingen, met de kleinste-kwadratenrechte $y = 0{,}6x + 2{,}2$. Na drie uur stijgt de score nauwelijks nog: het lineaire model past maar matig ($r \approx 0{,}77$).}}

\begin{werkwijze}
Controleer bij elke figuur: assen met naam en eenheid, een duidelijke schaal, metingen als punten (en een model als lijn), een nummer en een bijschrift dat zegt wat je moet zien. Verwijs in de tekst naar de figuur met haar nummer.
\end{werkwijze}

\begin{oefeningen}
\oef{h3-figuren}

\needspace{12\baselineskip}
\oef{h3-bijschrift}
\end{oefeningen}

\subsection{Tabellen}

Een figuur toont een verloop, een tabel toont exacte waarden. Ook een tabel moet zonder de tekst te lezen zijn. De grootheid en haar eenheid staan in de kop, en niet bij elk getal. Alle getallen van een meting hebben evenveel decimalen, zodat je ziet hoe nauwkeurig er gemeten is. Boven de tabel staat een nummer en een titel.

\voorbeeld{}{Verbeter deze tabel met metingen aan een slinger.\label{vb:h3-tabel}}{%
\uitleg{We zetten de eenheden in de kop, en schrijven alle lengtes in dezelfde eenheid en met evenveel decimalen.}
\figuren{{\renewcommand{\arraystretch}{1.2}%
\begin{tabular}{|c|c|c|c|}
\hline
lengte & 20 cm & 0,4 m & 60cm \\ \hline
tijd & 9,0 s & 12,7 s & 15,5 \\ \hline
\end{tabular}}\hspace{1cm}$\longrightarrow$\hspace{1cm}%
\begin{tabular}{@{}l@{}}
\textit{Tabel 1: tijd voor tien periodes.}\\[1mm]
{\renewcommand{\arraystretch}{1.2}%
\begin{tabular}{|c|c|c|c|}
\hline
$L$ (m) & $0{,}20$ & $0{,}40$ & $0{,}60$ \\ \hline
$10\,T$ (s) & $9{,}0$ & $12{,}7$ & $15{,}5$ \\ \hline
\end{tabular}}
\end{tabular}}
Links staan twee eenheden door elkaar, en bij $15{,}5$ ontbreekt de eenheid. Rechts staan de symbolen uit de tekst in de kop, met hun eenheid, en alle lengtes met twee decimalen: $0{,}4$ m zou doen denken dat die lengte minder nauwkeurig gemeten is.}

\begin{oefeningen}
\oef{h3-tabel}
\end{oefeningen}

\subsection{Presenteren}

Bij een presentatie kan het publiek niet terugbladeren. Het hoort alles maar één keer, en het onthoudt maar weinig. Daarom kies je vooraf de paar dingen die het zeker moet onthouden. Elke dia heeft één boodschap, en de titel zegt die boodschap. Een figuur werkt beter dan tekst, en een paar getallen beter dan een volle tabel.

\voorbeeld{}{Ontwerp drie dia's om het slingeronderzoek in drie minuten voor te stellen.\label{vb:h3-presentatie}}{%
\uitleg{Het publiek moet de vraag, de aanpak en het antwoord onthouden: één dia voor elk. De titel van elke dia is een volzin.}
\begin{tabular}{@{}>{\RaggedRight\arraybackslash}p{1.1cm}@{\hspace{3mm}}>{\RaggedRight\itshape\arraybackslash}p{6.3cm}@{\hspace{4mm}}>{\RaggedRight\small\itshape\arraybackslash}p{6.9cm}@{}}
\textbf{dia} & \textup{\textbf{titel}} & \textup{\textbf{inhoud}} \\[.4ex]
1 & Hangt de periode van een slinger af van zijn lengte? & Een foto van de opstelling, en de vraag in grote letters. \\[.6ex]
2 & We maten tien periodes, voor vijf lengtes. & Een schets met $L$ en $T$, en waarom tien periodes nauwkeuriger is. \\[.6ex]
3 & $T^2$ is evenredig met $L$. & De rechte van $T^2$ tegen $L$, en het besluit $T \approx 2{,}01\sqrt{L}$ in één regel.
\end{tabular}}

\begin{oefeningen}
\oef{h3-dia}
\end{oefeningen}

\opdrachtenkop

\begin{oefeningen}
\needspace{14\baselineskip}
\oef{h3-o6}

\oef{h3-o1}

\oef{h3-o7}

\oef{h3-o8}

\oef{h3-o9}

\needspace{12\baselineskip}
\oef{h3-o10}

\oef{h3-o2}

\oef{h3-o14}

\needspace{14\baselineskip}
\oef{h3-o11}

\oef{h3-o13}

\oef{h3-o15}

\oef{h3-o12}

\oef{h3-o4}

\oef{h3-o3}

\oef{h3-o17}

\needspace{14\baselineskip}
\oef{h3-o16}

\needspace{14\baselineskip}
\oef{h3-o18}

\oef{h3-o5}
\end{oefeningen}

%----------------------------------------------------------------------------------------
% KOPPEL: verbanden met andere leerstof
%----------------------------------------------------------------------------------------
\needspace{20\baselineskip}
\opdrachtsectie{Koppel}

\lkr{Logica en bewijzen en Statistiek: kijk na wat al gezien is. Het pakket legt de volgorde niet vast.}

\begin{oefeningen}
\oef{h3-koppel-1}

\oef{h3-koppel-2}
\end{oefeningen}

%----------------------------------------------------------------------------------------
\opdrachtsectie{Test jezelf}

\begin{oefeningen}
\oef{h3-t1}

\oef{h3-t2}

\oef{h3-t3}

\oef{h3-t4}

\oef{h3-t5}
\end{oefeningen}

\herhaalschema{%
  \fv{oef:h3-t1}{vb:h3-tekst}{oef:h3-schrijven}%
  \fv{oef:h3-t2}{vb:h3-figuur}{oef:h3-figuren}%
  \fv{oef:h3-t3}{vb:h3-bron}{oef:h3-bronnen}%
  \fv{oef:h3-t4}{vb:h3-besluit}{oef:h3-besluit}%
  \fv{oef:h3-t5}{vb:h3-nauwkeurig}{oef:h3-afronden}}

\ifzelfstudie
\section*{Maak je eigen samenvatting}
Sluit dit boek en schrijf in je schrift: de onderdelen van een verslag en wat erin hoort, het verschil tussen besluit en reflectie, wat je bij een bron vermeldt, de regels voor wiskundig schrijven, absolute en relatieve fout met de twee eigenschappen over fouten (en hun bewijs), de vuistregels voor afronden, en waaraan een goede figuur, tabel en dia voldoen. Maak er een checklist van die je bij je eigen verslag gebruikt.
\fi
